\chapter{推荐展示}
\label{chap:rs-presentation}

排序器给三篇登山包测评都打了高分，首页是否就该连续展示这三篇？用户可能只需要一篇测评，还想知道雨天如何保养背包。把两条单独表现不错的广告挨在一起，也可能使用户直接离开。即使对象完全相同，首屏的位置、摘要里的信息以及今天已经出现过的次数，仍会改变这次展示的结果。

上一章回答了怎样比较对象，本章把决策单位扩大为完整展示方案。内容、位置、表达、时间和交互共同决定用户实际遇到什么。各维度可以分别控制，也可以由同一个求解器联合确定；讨论时须说明哪些变量已经固定。从已有候选组合到目录标识生成，各种方法决定了对象怎样组成列表；表达、重复与用户控制则进一步改变展示体验。这些选择都需要利用反馈学习，并检验整组效果。多目标的取舍与跨请求的后果分别留给后两章。

% Outline: 5.1
\section{展示任务与效用}
\label{sec:rs-presentation-task}

设请求信息为$\bm{x}_t$，可用目录为$I_t$。一次展示行动可以写成
\begin{equation}
 a_t=(S_t,\rho_t,\bm e_t,\tau_t,\kappa_t),
 \label{eq:rs-presentation-action}
\end{equation}
其中$S_t\subseteq I_t$为展示对象集合，$\rho_t$把对象分配到位置，$\bm e_t$表示标题、摘要等信息版本，$\tau_t$规定时机与频次控制，$\kappa_t$规定交互入口及响应规则。对于简单列表，内容与位置可合写为$L_t=(i_1,\ldots,i_k)$；对于页面，还要给出模块、行列和首屏范围。会话方案则包含后续收到反馈时的行动规则，不能仅用一个静态列表代替。

展示效用是执行方案后的结果价值。令$Y(a)$为在明确观察窗口内执行$a$所产生的结果向量，$g$把这些结果转成当前要比较的效用，则
\begin{equation}
 U(\bm x,a)=\mathbb E[g(Y(a))\mid\bm x],\qquad
 a^\star\in\argmax_{a\in A(\bm x)}U(\bm x,a).
 \label{eq:rs-presentation-utility}
\end{equation}
$A(\bm x)$是可执行行动集合，包含资格、数量及界面规则。这里的潜在结果写法定义了决策目标；从已有日志学出的$\widehat U$能否识别$U$，还要满足第~\ref{chap:rs-evaluation}~章的数据条件。给模型输出取名“效用”，不会使它自动成为真实效果。

整组效用未必等于对象分数之和。若两个测评完全重复，第二篇的额外价值可能很小；背包与防雨罩却可能互相补充。用户注意也有限：下屏对象可能没有进入视野，点击第一篇之后也可能不再查看其余内容。一个包含邻接项的教学表达是
\[
 \widehat U(\bm x,L)
 =\sum_{j=1}^{k}v_j(\bm x,i_j)
  +\sum_{j=1}^{k-1}h(\bm x,i_j,i_{j+1}),
\]
其中$v_j$已包含指定位置的价值，$h$表示额外的相邻影响。这个简化式用于展示单项与组合的区别；真实影响可能跨越多个位置，也可能通过布局、摘要和重复状态产生。

本章的贯穿页面包含背包测评、通勤收纳指南等自然内容，以及背包和水壶广告。已有候选方案读取$C_t\subseteq I_t$、分数和关系信息，再从中组合；直接生成方案可以读取历史与目录编码，产生物品标识，不必先接收外部候选分数。两者都要落到可执行的对象、位置及展示版本。图~\ref{fig:rs-presentation-dimensions}列出同一页面的五组变量；这些是并列的控制对象，并非五个必须依次运行的模块。

\begin{figure*}[htbp]
 \centering
 % 教学页面：五组并列控制变量；连线是作用关系，不表示执行顺序。
\begin{tikzpicture}[
 x=1mm,y=1mm,
 font=\figurefont\small,
 dim/.style={draw=BookBlue,fill=BookBlue!6,line width=0.85pt,
             text width=38mm,align=left,inner sep=3mm},
 card/.style={draw=BookRule,fill=white,line width=0.85pt},
 relation/.style={draw=BookMuted,line width=0.85pt,rounded corners=1.2mm}
]
 \path[use as bounding box] (0,0) rectangle (169,82);
 \draw[draw=BookInk,fill=BookPaper,line width=1.1pt]
      (57,1) rectangle (110,81);
 \node[anchor=west,font=\figurefont\small\bfseries,text=BookInk]
      at (61,76) {通勤与户外};
 \node[anchor=west,font=\figurefont\footnotesize,text=BookMuted]
      at (61,69) {09:40 · 广告已见1次};

 \draw[card] (61,49) rectangle (106,64);
 \node[anchor=west,text=BookInk] at (64,59) {$i_1$　背包测评};
 \node[anchor=west,font=\figurefont\footnotesize,text=BookMuted]
      at (64,53) {重量、收纳与使用条件};

 \draw[draw=BookAmber,fill=BookAmber!7,line width=0.85pt]
      (61,25) rectangle (106,45);
 \node[anchor=west,text=BookInk] at (64,40) {$i_2$　背包商品};
 \node[anchor=east,font=\figurefont\footnotesize,text=BookAmber]
      at (104,40) {广告};
 \node[anchor=west,font=\figurefont\footnotesize,text=BookInk]
      at (64,33) {399元 · 面料防泼水};
 \node[anchor=west,font=\figurefont\footnotesize,text=BookMuted]
      at (64,28) {不适合长时间淋雨};

 \draw[BookMuted,dashed,line width=0.85pt]
      (58,47) -- (109,47);
 \draw[card] (61,5) rectangle (106,21);
 \node[anchor=west,text=BookInk] at (64,16) {$i_3$　保养指南};
 \node[anchor=east,font=\figurefont\footnotesize,text=BookTeal]
      at (103,9) {收藏　不感兴趣};

 \node[dim] (content) at (23,65) {\textbf{内容}\\选入哪些对象、各来源数量};
 \node[dim] (position) at (23,43) {\textbf{位置}\\顺序、模块与首屏范围};
 \node[dim] (expression) at (23,17) {\textbf{表达}\\摘要、条件与商业标识};
 \node[dim] (timing) at (146,68) {\textbf{时间}\\本次时机、间隔与累计频次};
 \node[dim] (interaction) at (146,10) {\textbf{交互}\\可用入口与操作后的响应};

 \draw[relation] (content.east) -- (52,65) -- (52,59) -- (61,59);
 \draw[relation] (position.east) -- (52,43) -- (52,47) -- (58,47);
 \draw[relation] (expression.east) -- (52,17) -- (52,33) -- (61,33);
 \draw[relation] (timing.west) -- (110,68);
 \draw[relation] (interaction.west) -- (110,10);
\end{tikzpicture}

 \caption{同一推荐页面的五个展示维度。页面对象为教学示意；连线把每组决策关联到其控制的页面元素或状态，不表示算法执行顺序。广告数量、版位和表达可以联合改变，右侧时间及交互规则还会影响后续请求。}
 \label{fig:rs-presentation-dimensions}
\end{figure*}

% Outline: 5.2
\section{展示方案设计}
\label{sec:rs-presentation-design}

设计方案时，先圈定可改动的变量，再讨论如何求解。只接收固定广告版位的算法，不能因为使用了“整页”响应预测就被解释为已经优化广告负载。相反，直接输出完整有序列表的模型已经确定了部分位置，不应再把这些位置当作无条件自由变量。下面按五个维度展开，每种方法都保留其输入和控制范围。

% Outline: 5.2.1
\subsection{展示内容组织}
\label{sec:rs-content-organization}

内容组织决定展示哪些对象、各来源各占多少，以及它们共同满足什么需求。去重排除同一对象或等价副本；多样性可以要求主题、作者或表达角度的差异；方面覆盖则要求照顾当次需求中的多个部分。它们并不等价：推荐许多不同品牌的登山包仍可能只回答“买什么包”，没有回答防水、收纳和维护问题。

自然内容与广告合流时，还要比较跨来源的价值。第~\ref{sec:rs-score-calibration}~节的概率校准使同一响应口径下的分数可比较，但观看概率、广告收入和关闭成本仍有不同单位。将它们组合必须声明换算系数或约束。广告支付又涉及其他参与者，不能仅凭混排分数决定收费，相关机制见第~\ref{chap:rs-advertising}~章。

\begin{table*}[htbp]
 \centering\small
 \begin{tabularx}{\linewidth}{@{}p{28mm}XX@{}}
  \toprule
  方法 & 主要输入 & 直接输出与固定条件\\
  \midrule
  MMR、DPP & 候选、质量与相似关系 & 选择子集；位置是否沿用构建顺序另定\\
  Ads Allocation & 两路已排序内容、响应与收入估计 & 合法混排；保留各来源内部顺序\\
  Virtual Bids & 固定自然内容、广告版位与候选广告 & 广告身份及次序；数量和版位固定\\
  Ad-load Balancing & 请求上下文、可选负载、随机日志 & 一次获取的负载；再落实广告对象\\
  PRM、Seq2Slate、NAR4Rec & 既有候选及其特征 & 候选内的次序或展示列表\\
  OneRec & 历史语义标识、目录编码、列表反馈 & 从目录代码生成多物品序列\\
  OneRec-Think & 历史标识、文本与当次要求 & 理由和下一物品或生成候选；仍需组织展示\\
  \bottomrule
 \end{tabularx}
 \caption{内容组织方法的接口与控制范围。比较时应固定实际可用信息、目录范围、展示规则和计算预算；方法名称本身不决定它是否使用外部候选、是否已经安排位置。}
 \label{tab:rs-content-interfaces}
\end{table*}

表~\ref{tab:rs-content-interfaces}中的方法可以组合。例如负载策略先确定广告数量，广告分配器再选择对象；生成的候选也可进入去重与校准。组合时需要核验接口：第~\ref{sec:rs-onerec-think}~节的OneRec-Think以物品生成及文字理由为主，第~\ref{sec:rs-onerec-v2}~节的OneRec-V2主体也是单目标物品，二者都不能仅因名称相近就沿用OneRec的会话列表设定。

% Outline: 5.2.1.1
\subsubsection{MMR：相关性与冗余重排}
\label{sec:rs-mmr}

已有分数较高的候选可能互相重复。\term{最大边际相关性}（maximal marginal relevance, MMR）在逐步构建集合时，一面奖励与需求的相关性，一面惩罚与已选对象的最大相似度\scite{rs-h12}设$r_i$为相关性，$s_{i,j}$为对象相似度，$S$为已选集合，下一项为
\begin{equation}
 i^\star=\argmax_{i\in C\setminus S}
 \left\{\lambda r_i-(1-\lambda)\max_{j\in S}s_{i,j}\right\}.
 \label{eq:rs-mmr}
\end{equation}
空集合时把冗余项设为零，选入$i^\star$后更新$S$，直到目标长度。相关性和相似度可以由前章的模型提供，MMR本身不训练这两种估计；$\lambda\in[0,1]$应按明确的验证目标选择，分数尺度也必须先约定。

用三条教学候选计算：$a$和$b$是同类背包测评，$c$是雨天保养指南。令$(r_a,r_b,r_c)=(0.90,0.85,0.70)$，$(s_{a,b},s_{a,c},s_{b,c})=(0.95,0.10,0.20)$，相似关系对称，取$\lambda=0.7$。第一步三个值为0.630、0.595和0.490，选$a$。第二步$b$的值为$0.595-0.3\times0.95=0.310$，$c$的值为$0.490-0.3\times0.10=0.460$，故选$c$。基础相关性较低的对象通过较少冗余进入结果。

输出既有集合$\{a,c\}$，也有构建次序$(a,c)$；是否直接将该次序用于页面，要看位置效用。最大相似度只惩罚某一种重复，不保证所有需求方面都有代表，也不保证某个任意整组目标的全局最优。若两件相似物品恰好便于价格比较，过强惩罚还可能破坏当前任务。

% Outline: 5.2.1.2
\subsubsection{DPP：质量与差异的集合选择}
\label{sec:rs-dpp}

MMR只查看最相似的已选对象，无法完整表达一个候选能否被多个已选对象共同替代。\term{行列式点过程}（determinantal point process, DPP）用子集对应向量张成的体积表达差异；向量方向近似重合时，体积减小。给候选$i$一个非负质量$q_i$及单位特征$\bm z_i$，定义半正定核
\begin{equation}
 K_{i,j}=q_iq_j\langle\bm z_i,\bm z_j\rangle,\qquad
 \mathbb P(S)=\frac{\det(\bm K_{S,S})}{\det(\bm I+\bm K)}.
 \label{eq:rs-dpp}
\end{equation}
$\bm K$是候选上的核矩阵，$\bm I$为同阶单位矩阵。这个概率模型给每个子集分配概率；寻找概率最大的子集是最大后验（MAP）推断，并不等于从DPP采样。固定长度$k$时，比较所有$|S|=k$的子集，其归一化常数虽改变，最大化行列式的目标不变。

直接为每个扩展子集重新分解矩阵代价很高。快速贪心推断维护Cholesky投影和残差\scite{rs-r15}设当前选中子核为$\bm K_{S,S}=\bm V\bm V^\top$，未选对象$i$的投影$\bm c_i$满足$\bm V\bm c_i=\bm K_{S,i}$，残差平方为$d_i^2=K_{i,i}-\lVert\bm c_i\rVert_2^2$。块行列式给出
\[
 \det(\bm K_{S\cup\{i\},S\cup\{i\}})
 =\det(\bm K_{S,S})d_i^2.
\]
因而每步选择$\log d_i^2$最大的剩余项。若本步选$j$，对每个剩余$i$只需计算
\begin{equation}
 e_i=\frac{K_{j,i}-\langle\bm c_j,\bm c_i\rangle}{d_j},
 \quad
 \bm c_i\leftarrow[\bm c_i;e_i],
 \quad
 d_i^2\leftarrow d_i^2-e_i^2.
 \label{eq:rs-dpp-update}
\end{equation}
在给定核或能够按需取核元素的条件下，选取$k$项的这部分计算为$O(|C|k^2)$，还应另计特征及核的取得成本。

沿用$a,b,c$，这次为单独观察差异而把三个质量都设为1。取$K_{a,b}=0.9$、$K_{a,c}=0.2$、$K_{b,c}=0.3$，对角线均为1；该教学核正定。若并列时先选$a$，则$b$残差为$1-0.9^2=0.19$，$c$为$1-0.2^2=0.96$，分别也是两个二阶子核的行列式。贪心选择$c$。再更新$b$时，新投影为$(0.3-0.9\times0.2)/\sqrt{0.96}\approx0.1225$，残差降至0.175。

固定长度时可能需要接受负的对数增益；无长度约束的贪心则可在最大增益不再为正时停止。残差近零意味着几乎线性相关，不能继续除以零；核秩不足时也未必存在正概率的指定长度集合。数值截断应区分浮点误差与实质退化。快速实现加速的是贪心规则，没有把一般DPP的MAP问题变成全局精确求解；多样性含义仍由$\bm z_i$决定。

% Outline: 5.2.1.3
\subsubsection{Ads Allocation：有序来源的约束混排}
\label{sec:rs-ads-allocation}

自然内容和广告已各自完成排序时，重新打乱两路内部次序会破坏上游的价值判断，也可能影响已经关联的竞价规则。\term{Ads Allocation in Feed}将混排放在来源排序之后，研究保留内部顺序时的广告插入\scite{rs-m01}设$\mathcal L_{\mathrm{merge}}$为合法合流列表，$\widehat R(L)$与$\widehat E(L)$分别为收入和参与度估计，一种任务写法是
\begin{equation}
 \max_{L\in\mathcal L_{\mathrm{merge}}}\widehat R(L)
 \quad\text{s.t.}\quad \widehat E(L)\geq e_{\min}.
 \label{eq:rs-ads-allocation}
\end{equation}
广告数量、间隔及禁止首位等规则进入可行集合。响应来自已有模型，阈值来自目标要求；论文还讨论通过约束的对偶变量得到效用换算系数，以及不同位置和间隔假设下的轻量算法。

为了看清可行集合，取自然队列$(o_1,o_2)$、广告队列$(a_1,a_2)$，四项全部展示，要求两广告不相邻。每步只能取某一路尚未使用的首项。六种保序合流中只有表~\ref{tab:rs-merge-example}的三种合法。教学假设位置折扣为$(1,0.8,0.6,0.4)$，基础参与概率按$(o_1,o_2,a_1,a_2)$为$(0.6,0.5,0.2,0.1)$，两个广告的基础预期收入为1.5元与1元；位置折扣同时乘到对应参与和收入上，暂不添加其他组合影响。

\begin{table}[htbp]
 \centering\small
 \begin{tabular}{@{}lrrc@{}}
  \toprule
  合法列表 & 参与度 & 收入／元 & 参与度$\geq1$\\
  \midrule
  $(a_1,o_1,a_2,o_2)$ & 0.94 & 2.10 & 否\\
  $(a_1,o_1,o_2,a_2)$ & 1.02 & 1.90 & 是\\
  $(o_1,a_1,o_2,a_2)$ & 1.10 & 1.60 & 是\\
  \bottomrule
 \end{tabular}
 \caption{保留两路内部顺序的教学混排。参与度是四个折扣后参与概率之和，表示预期响应次数；收入为给定估计，不从广告出价直接等同得到。}
 \label{tab:rs-merge-example}
\end{table}

例如第二行的参与度为$0.2+0.8\times0.6+0.6\times0.5+0.4\times0.1=1.02$。阈值为1时，第一行虽收入最高却不可行，最终选择第二行。这个枚举给出了对象身份、插入位置与混合顺序，之后页面只需执行已确定的分配。

四对象枚举用于核验任务表述，不支持工业规模的效率结论。实际问题需要利用位置响应、间隔影响或可分解性设计求解器；如果这些假设改变，原有最优性结论也要重新检查。要求保留来源顺序的算法更不能直接被当成任意候选的整页重排器。

% Outline: 5.2.1.4
\subsubsection{Virtual Bids：固定版位下的广告组合}
\label{sec:rs-virtual-bids}

相邻广告会竞争或互相补充，广告还会与自然内容共同影响点击。\term{Virtual Bids}的原应用固定自然商品及其位置，也固定广告数量和版位，仅决定哪些广告入选及它们之间的次序\scite{rs-m02}整组响应模型读取用户、广告、自然商品及位置，输出各广告在方案$L$下的点击概率$\widehat p_j(L)$。用$b_j$表示相应广告的每点击出价，$\nu$表示平台对一次广告点击的虚拟估值，其分配目标可写成
\begin{equation}
 F_\nu(L)
 =\underbrace{\sum_{j\in P_{\mathrm{ad}}}b_j\widehat p_j(L)}
              _{\text{出价加权响应项}}
  +\underbrace{\nu\sum_{j\in P_{\mathrm{ad}}}\widehat p_j(L)}
              _{\text{平台虚拟价值项}}.
 \label{eq:rs-virtual-bids}
\end{equation}
$P_{\mathrm{ad}}$为固定广告位置集合。$\nu$把参与度换算到货币量纲；这一项不是真实付款。第一项也是分配中的出价价值或收入代理，最终现金收入取决于支付机制，不能把出价直接记成收款。

取两个固定广告位置和三个广告$A,B,C$，出价分别为每点击3、2、1元。表~\ref{tab:rs-virtual-bids-example}给出教学列表模型对六种有序组合的点击估计，以及不计和计入虚拟价值时的分数。自然内容和广告版位在所有行中完全相同。

\begin{table}[htbp]
 \centering\small
 \begin{tabular}{@{}lrrrr@{}}
  \toprule
  广告组合 & $\widehat p_1$ & $\widehat p_2$ & $F_0$／元 & $F_1$／元\\
  \midrule
  $(A,B)$ & 0.14 & 0.08 & 0.58 & 0.80\\
  $(B,A)$ & 0.10 & 0.10 & 0.50 & 0.70\\
  $(A,C)$ & 0.10 & 0.22 & 0.52 & 0.84\\
  $(C,A)$ & 0.18 & 0.09 & 0.45 & 0.72\\
  $(B,C)$ & 0.10 & 0.16 & 0.36 & 0.62\\
  $(C,B)$ & 0.17 & 0.09 & 0.35 & 0.61\\
  \bottomrule
 \end{tabular}
 \caption{固定广告版位下的教学组合。$F_0$只计算出价加权响应，$F_1$额外按每点击1元计入平台虚拟价值；两者均为分配分数，未计算最终支付。}
 \label{tab:rs-virtual-bids-example}
\end{table}

当$\nu=0$时选$(A,B)$。若平台每次广告点击的虚拟估值为1元，$(A,C)$以0.84元超过$(A,B)$的0.80元。变化来自整组响应与参与价值，不是简单地给每件广告加上同一个常数。即使自然对象没有改动，响应预测仍需看到它们，因为它们参与了用户选择。

虚拟估值不是点击模型的一个普通特征。论文通过历史请求上的外层搜索选择估值：内层在每次请求的候选组合中分配，外层比较得到的参与度与收入代理相对各自理想目标的距离。这给出特定取舍准则，不代表数据能够独自决定平台应重视什么。截取少量广告再枚举可降低成本，也会限制搜索空间。输出广告分配后，考虑外部性的支付计算由第~\ref{chap:rs-advertising}~章展开；本应用没有联合优化自然位置或广告负载。

% Outline: 5.2.1.5
\subsubsection{Ad-load Balancing：一次获取的负载学习}
\label{sec:rs-ad-load}

固定每页两个广告没有利用用户差异：有的人刚进入会话，有的人已经连续浏览多次。\term{Ad-load Balancing}把一次信息流获取的广告负载作为上下文行动，用随机日志学习或比较分配策略\scite{rs-m05}上下文包含用户、会话阶段及内容市场信息，行动可以是广告数量$n_{\mathrm{ad}}$，也可以是数量与位置的合法组合；反馈同时保留收益与滚动、退出等体验代理。

设日志为$(\bm x_t,a_t,r_t,\mu(a_t\mid\bm x_t))$，其中$\mu$为收集数据时的行动概率。目标策略$\pi$在同类请求中选择哪种负载，不能只比较日志里各行动的原始均值，因为它们出现的用户可能不同。离策略估计利用记录概率作分布修正，或结合结果预测降低方差；逆倾向评分与双重稳健估计的推导见第~\ref{chap:rs-evaluation}~章。

一个教学随机策略在同类页面上以$(0.2,0.5,0.3)$的概率展示零、一个或两个广告。假设某记录执行两个广告，成熟后的收益、退出标记为$(0.9\text{元},1)$；另一次执行一个广告，结果为$(0.5\text{元},0)$。如果教学效用取收益减去0.4元乘退出标记，两次观测都得到0.5元，却不能凭这两条记录认定两负载等效：还需要其他样本、上下文和概率修正。若某用户组经充分数据估计的三种效用为$(0,0.4,0.5)$，本轮选择两个；另一组为$(0,0.3,0.1)$，选择一个。

负载确定后仍要选择具体广告，并检查版位与间隔。行动只描述数量时，不同的广告选择器会改变该数量对应的结果分布，所以训练与使用应保持这一执行接口明确。日志从未尝试的负载没有可靠支撑；预测器的外推不等于已经观察过。一次获取的上下文收益，也不等于第~\ref{chap:rs-long-term}~章显式考虑未来状态的连续决策目标。

% Outline: 5.2.1.6
\subsubsection{OneRec：列表生成与偏好对齐}
\label{sec:rs-onerec}

前面的方法都在已有候选或来源流上做决定。\term{OneRec}将用户历史编码后，直接生成多个物品的语义标识，把列表内部的选择关系纳入同一个自回归模型\scite{rs-r08}语义标识的构造沿用第~\ref{sec:rs-tiger}~节的目录编码思想；模型的编码器读取历史代码，解码器读取历史表示和已经生成的目标代码。稀疏专家网络用于扩大模型容量，会话列表而非单个下一物品作为训练目标。

设输出列表$L=(i_1,\ldots,i_k)$的代码串为$(z_1,\ldots,z_m)$，其中含物品边界，$m$为总代码长度。对筛选出的高质量会话，下一代码预测损失为
\begin{equation}
 \mathcal L_{\mathrm{NTP}}
 =-\sum_{\ell=1}^{m}\log
   p_\theta(z_\ell\mid H_{u,t},z_{<\ell}).
 \label{eq:rs-onerec-ntp}
\end{equation}
后一个物品的概率依赖前面已生成的物品，因此可以学习列表关系。训练时输入真实前缀，推理时使用自身前缀，仍存在误差累积；会话筛选规则也决定了模型模仿哪种行为。

同一用户通常只看到一次实际列表，无法同时为多个生成列表提供真实优劣反馈。OneRec因此训练会话奖励模型：读取用户与列表的联合表示，利用历史会话的观看和交互监督多个结果头，再给模型产生的候选列表评分。迭代偏好对齐（iterative preference alignment, IPA）从束搜索结果中挑选奖励较高和较低的列表，构成$(L^+,L^-)$。直接偏好优化（DPO）的通用推导见\BookSectionRef{sec:alignment-learning}{基础卷的“训练阶段的数据与参数干预”一节}；在这里，更新对象是生成列表的策略，而偏好标记主要来自奖励模型。

为看清实际更新量，记固定参考策略为$p_{\mathrm{ref}}$，本轮列表对的相对对数概率差为
\begin{align}
 d_\theta&=
 \log\frac{p_\theta(L^+\mid H)}{p_{\mathrm{ref}}(L^+\mid H)}
 -\log\frac{p_\theta(L^-\mid H)}{p_{\mathrm{ref}}(L^-\mid H)},\\
 \mathcal L_{\mathrm{align}}
 &=\mathcal L_{\mathrm{NTP}}
   -\lambda\log\sigma(\beta d_\theta).
 \label{eq:rs-onerec-dpo}
\end{align}
$\beta>0$控制偏好损失的尺度，$\lambda$控制两项训练贡献。训练从当前生成器取得新列表对，更新生成器后再刷新样本；并非永久反复拟合同一组固定列表。奖励模型对未真实曝光列表的评分仍是估计，不是用户已作出的选择。

教学例中，三条历史分别为背包评测、雨天徒步和通勤收纳，代码依次是$(1,2)$、$(3,1)$、$(2,4)$。解码器生成物品边界后的$(1,5)$和$(3,2)$，目录映射到一篇新测评和一篇防雨保养指南。执行前核验两代码是否对应现存物品、是否重复、是否仍具展示资格；若第二件已下架，应按声明的补充规则处理，而不能把未知代码直接渲染为商品。若候选列表$L^+$为这两项，$L^-$为两篇重复测评，奖励模型可以产生偏好对，却仍需单独检验其评分是否可信。

假设参考策略对两列表的概率相同，当前概率比为2，取$\beta=1$，则该对的DPO损失为$-\log\sigma(\log2)=-\log(2/3)\approx0.405$。概率比增大将降低这项损失，但不能保证列表合法，更不能保证在线体验改善。本节依据2025年公开预印本介绍机制；论文的平台试验有其数据和服务范围。比较多阶段方案时，还需同时报告实际可访问目录、解码预算、资格处理以及独立在线证据。

% Outline: 5.2.1.7
\subsubsection{PRM：列表上下文的个性化重排}
\label{sec:rs-prm}

如果上游已经给出了不错的候选，仍可以让每个对象的分数看到同组其他对象。\term{个性化重排模型}（personalized re-ranking model, PRM）把候选特征、个性化向量和原次序的位置嵌入交给Transformer编码，输出各候选的列表相关分数\scite{rs-r44}其中个性化向量可取自预训练响应网络靠近输出层的表示，因此同一物品面向不同用户可以有不同编码。

将候选$i$的原始特征$\bm x_i$与个性化表示$\bm v_{u,i}$拼接，再加入对应位置嵌入并投影，得到输入行$\bm e_i$。自注意力在所有候选行之间交换信息，产生$\bm h_i$；原文用线性头和列表softmax给出
\[
 q_i=\frac{\exp(\bm w^\top\bm h_i+b)}
 {\sum_{j\in C}\exp(\bm w^\top\bm h_j+b)},
 \qquad
 \mathcal L=-\sum_{i\in C}y_i\log q_i.
\]
$y_i$为该列表中的点击标签。这个$q_i$是候选范围内归一化的重排分数，不应不经验证就解释为每件物品独立的曝光点击率；同一列表可能有多次点击。

考虑一个单头、标量值的教学注意力计算。对候选$a$，若它给$(a,b,c)$的注意权重为$(0.2,0.6,0.2)$，三个值为$(1,0.8,0.1)$，则聚合值为0.70。把$b$换成值为$-0.4$的另一候选，并暂固定权重以隔离值通路，聚合值变成$0.2-0.24+0.02=-0.02$。真实模型的键和权重也会改变，但这个局部计算已说明：$a$自身未变，上下文仍能改变其输出；逐对象评分不具备这条输入路径。

服务时按输出分数重排即可。标准全注意力对$n$个候选的关系计算随$n^2$增长，适合在已缩小的范围内使用。更关键的是，预测所见的初始排列与最终执行排列可能不同；上下文评分并未自动搜索所有排列。多样性、广告间隔和资格仍需检查，训练日志对哪些排列有支撑也需明确。

% Outline: 5.2.1.8
\subsubsection{Seq2Slate：指针解码的列表选择}
\label{sec:rs-seq2slate}

PRM一次输出全部分数，后续对象不能在选择时重新读取“已经选了什么”。\term{Seq2Slate}用LSTM编码候选，再由另一个LSTM逐步选择下一个候选；指针式解码表示输出指向当前输入集合中的对象，而不是固定全目录词表\scite{rs-r45}候选编码为$\bm e_i$、第$j$步解码状态为$\bm d_j$，指向$i$的logit为
\begin{equation}
 s_{j,i}=\bm v^\top\tanh(\bm W_{\mathrm{enc}}\bm e_i+
                          \bm W_{\mathrm{dec}}\bm d_j),
 \quad
 p_j(i)=\frac{\exp(s_{j,i})}
 {\sum_{\ell\in C\setminus S_{j-1}}\exp(s_{j,\ell})}.
 \label{eq:rs-seq2slate}
\end{equation}
该概率只对$i\notin S_{j-1}$定义，已选对象概率设为零。选中对象的表示进入下一解码步，故整个列表的概率为$\prod_j p_j(i_j)$，重复屏蔽与前缀条件各有作用。

设四个教学候选第一步的指数logit为$(4,3,2,1)$。归一化后，四项概率依次为0.4、0.3、0.2和0.1，贪心选择$a$。第二步在读入$a$后，四个指数logit变为$(5,2,4,1)$。将$a$屏蔽后，$(b,c,d)$概率为$(2/7,4/7,1/7)$，选择$c$，两步路径概率为$0.4\times4/7\approx0.229$。若不屏蔽，分数最高的$a$会再次入选；若仅屏蔽但不更新解码状态，也没有真正利用前缀关系。

点击日志只提供哪些对象被点，不提供唯一正确排列。原文给出两类学习方法：用点击标签计算候选排列的排名奖励，通过REINFORCE提高较好排列的概率；或在每一步对剩余候选应用点击监督损失，并移除已选对象的标签，按位置权重累计。前者可用$(R-b)\nabla_\theta\log p_\theta(L)$估计梯度，$b$为降低方差的基线；后者不需要把所有点击对象人为排成唯一答案。

使用历史点击评价新排列时，仍隐含标签能否迁移到新位置与新组合的问题，不能把弱监督训练自动视为无偏的策略价值学习。选择$k$项需要$k$次解码，各步在$n$个候选上比较，固定隐藏维度时为$O(nk)$量级，另计编码成本。Seq2Slate始终受输入候选覆盖限制，与OneRec从目录代码生成列表的接口不同。

% Outline: 5.2.1.9
\subsubsection{NAR4Rec：并行列表生成与评价}
\label{sec:rs-nar4rec}

逐项运行大模型会增加列表生成延迟，但完全独立选每个位置又容易产生重复或不协调。\term{NAR4Rec}先并行计算位置与候选的匹配，再用依赖已选对象的对比解码构造方案，最后交给列表评价器选择\scite{rs-r75}候选编码器以自注意力处理对象，位置编码器以位置嵌入、自注意力和对候选的交叉注意力产生位置表示。

设候选表示为$\bm h_i$，位置表示为$\bm z_j$。匹配概率
\[
 P_{j,i}=\frac{\exp(\langle\bm z_j,\bm h_i\rangle)}
                  {\sum_{\ell\in C}\exp(\langle\bm z_j,\bm h_\ell\rangle)}
\]
可以对所有$(j,i)$并行计算。原文依据整组反馈效用是否超过阈值区分训练列表：较高效用列表使用$-\sum_j\log P_{j,i_j}$，较低效用列表使用$-\sum_j\log(1-P_{j,i_j})$。后者压低这些位置—对象匹配的概率，是逐位置的不似然损失，不能写成仅对整列表概率使用一次$-\log(1-p(L))$。额外的表示对比项帮助区分不同物品和位置。

解码时，在未选候选中比较
\[
 (1-\alpha)P_{j,i}
 -\alpha\max_{\ell\in S_{j-1}}
       \operatorname{cos}(\bm h_i,\bm h_\ell),
\]
空前缀时相似度惩罚为零。其形式与MMR相近，但相关性项来自已经并行计算的位置匹配矩阵。生成多个候选排列后，评价器重新编码每个完整排列，预测其中各对象的反馈并聚合效用；执行估计总效用较高的合法方案。

取四候选、两个位置的教学概率矩阵
\[
 \bm P=
 \begin{pmatrix}
  0.45&0.30&0.15&0.10\\
  0.40&0.35&0.15&0.10
 \end{pmatrix}.
\]
独立取每行最大值会选两次$a$。若第一位选$a$，并取$\alpha=0.4$、$a$与$(b,c,d)$的相似度为$(0.9,0.1,0)$，第二步三个值为$-0.15,0.05,0.06$，于是选$d$。若另一生成配置产生$(a,c)$，评价器给$(a,d)$、$(a,c)$的教学效用分别为0.52和0.58，则选后者。匹配概率最高与整组评价最高不是同一标准。

“非自回归”描述了概率矩阵的计算方式，最终对比解码仍逐步读取已选集合；不能据此声称整个推理没有序列依赖。多个方案的评价也要计入总延迟。模型只能重排已有候选，评价器误差和训练列表的选择范围仍限制结果，较高的预测效用需要在线反馈验证。

% Outline: 5.2.1.10
\subsubsection{Calibrated Recommendations：兴趣构成的贪心重排}
\label{sec:rs-calibrated-recommendations}

差异更大的列表未必符合用户兴趣。若用户历史中80\%是爱情片、20\%是动作片，把五个名额分给五种从未表达过兴趣的类型，会提高类别数，却未必照顾这两类需求。\term{列表校准}（calibrated recommendations）要求展示的类别构成接近一个明确的兴趣分布；这里校准的是内容比例，与前章pCTR校准不同\scite{rs-r110}

设$G$为类别集合，$p_g$为依据历史窗口和权重估计的目标分布，$q_g(S)$为所选列表的类别分布。多标签对象可用归一化的类别贡献，且历史与列表采用一致口径。为防止列表缺少某类时散度无穷大，令
\begin{align}
 \widetilde q_g(S)&=(1-\epsilon)q_g(S)+\epsilon p_g,\\
 D(S)&=\sum_{g\in G}p_g\log\frac{p_g}{\widetilde q_g(S)},\\
 F(S)&=(1-\lambda)\sum_{i\in S}r_i-\lambda D(S).
 \label{eq:rs-calibrated-greedy}
\end{align}
$\epsilon>0$为小的平滑系数，$r_i$为基础相关性，$\lambda$控制取舍。贪心每步加入使$F(S\cup\{i\})$最大的对象，直至指定长度。已有模型提供$r_i$，该重排不另行训练CTR模型。

教学候选中五部爱情片$R_1,\ldots,R_5$的相关性为0.95、0.90、0.85、0.80、0.75，两部动作片$A_1,A_2$为0.70、0.65；类型互斥，列表内各位置等权。取$p=(0.8,0.2)$、$\epsilon=0.01$、$\lambda=0.5$。第一步只选爱情片时$D\approx0.7441$，只选动作片时$D\approx3.3639$，故选$R_1$。第二步若添$R_2$，$F\approx0.5529$；若添$A_1$，构成为$(0.5,0.5)$、$D\approx0.1892$，$F\approx0.7304$，故选$A_1$。

随后依次选$R_2,R_3,R_4$，四步和五步后的构成分别为$(0.75,0.25)$与$(0.8,0.2)$，最终散度为零。输出为四部爱情片、一部动作片；中间步骤偏向补齐次要兴趣，最终比例由整个长度共同决定。构建次序不必就是最后的页面次序，若使用位置权重，还要重新计算展示构成。

这只是给定数据和权重下的结果，短列表的整数粒度可能根本无法精确匹配比例，某类没有候选时也无法补齐。历史分布受旧曝光和选择影响，当次需求可能已经变化；复制历史比例没有天然的公平或满意度保证。应明确目标分布的证据，并在第~\ref{chap:rs-constraints}~章讨论其与其他目标的取舍。

% Outline: 5.2.1.11
\subsubsection{Minimum-Cost Flow Calibration：兴趣构成的流量求解}
\label{sec:rs-flow-calibration}

贪心一旦选入对象，通常不会再撤回早先决定。\term{最小费用流校准}将对象—槽位分配与类别计数放进同一个图，在明确的图目标上进行整体求解\scite{rs-r111}先考虑每件物品恰属一类、每个位置等量计入类别比例的情形。设$A_{i,j}$为将$i$放入槽位$j$的价值，$M_{i,j}\in\{0,1\}$为分配变量，$n_g$为类别$g$的入选数，列表长度为$k$。

该方法使用的构成惩罚是$D_{\mathrm{KL}}(q\Vert p)$，方向与式~\eqref{eq:rs-calibrated-greedy}的$D_{\mathrm{KL}}(p\Vert\widetilde q)$不同。不能把两者的算法成绩直接解释为只改变了求解器。定义单类计数代价
\[
 E_g(n)=\frac{n}{k}\log\frac{n}{kp_g},\qquad E_g(0)=0.
\]
将相关性收益取负，就得到需要最小化的总费用：
\begin{equation}
 \min_{\bm M}\left\{
 -(1-\lambda)\sum_{i,j}M_{i,j}A_{i,j}
 +\lambda\sum_g E_g(n_g)\right\}.
 \label{eq:rs-flow-objective}
\end{equation}
约束为每槽位恰好一个对象、每对象至多占一个槽位，且$n_g$与实际入选数一致。被允许选择的类别应有$p_g>0$，否则需事先规定平滑或禁止规则。

流网络从源点连到$k$个槽位，各边容量为1；槽位$j$连到候选$i$的入口，费用为$-(1-\lambda)A_{i,j}$。每件物品的入口与出口之间有容量为1的边，防止它被重复选择；出口连到所属类别节点。类别$g$再经单位容量计数边到汇点，第$n$条边的费用为$\lambda[E_g(n)-E_g(n-1)]$。这些增量单调不减，所以流会优先使用较小计数的边，累计费用正好等于$\lambda E_g(n_g)$。发送$k$单位流，就得到对象和槽位分配；增广过程中可以通过残余网络修订早先分配。

图~\ref{fig:rs-calibration-flow}只画出一条可用路径。所有槽位都能连接合法候选，但必须共享对应物品的容量边，因此不会把同一对象分给多个位置。类别后的计数通道承担构成代价，保持与对象相关性分开计费。

\begin{figure*}[htbp]
 \centering
 % 最小费用流的单条可用路径示意；全部边容量为1，未列出的费用为0。
\begin{tikzpicture}[
 x=1mm,y=1mm,font=\figurefont\small,
 vertex/.style={circle,draw=BookBlue,fill=BookBlue!6,
                line width=0.85pt,minimum size=8mm,inner sep=1mm},
 flow/.style={->,draw=BookInk,line width=1.1pt,>=latex}
]
 \path[use as bounding box] (0,0) rectangle (169,36);
 \node[vertex] (source) at (5,21) {$s$};
 \node[vertex] (slot) at (25,21) {$j$};
 \node[vertex] (itemin) at (59,21) {$i_{\mathrm{in}}$};
 \node[vertex] (itemout) at (84,21) {$i_{\mathrm{out}}$};
 \node[vertex] (category) at (109,21) {$g$};
 \node[vertex] (count) at (139,21) {$\ell$};
 \node[vertex] (sink) at (164,21) {$t$};
 \draw[flow] (source) -- (slot);
 \draw[flow] (slot) -- (itemin);
 \draw[flow] (itemin) -- (itemout);
 \draw[flow] (itemout) -- (category);
 \draw[flow] (category) -- (count);
 \draw[flow] (count) -- (sink);
 \node[text=BookMuted,font=\figurefont\footnotesize] at (5,32) {源点};
 \node[text=BookMuted,font=\figurefont\footnotesize] at (25,32) {槽位};
 \node[text=BookMuted,font=\figurefont\footnotesize] at (71.5,32) {物品入口与出口};
 \node[text=BookMuted,font=\figurefont\footnotesize] at (109,32) {类别};
 \node[text=BookMuted,font=\figurefont\footnotesize] at (139,32) {计数通道};
 \node[text=BookMuted,font=\figurefont\footnotesize] at (164,32) {汇点};
 \node[text=BookInk,font=\figurefont\footnotesize] at (42,9)
      {$-(1-\lambda)A_{i,j}$};
 \node[text=BookInk,font=\figurefont\footnotesize] at (71.5,9) {容量1};
 \node[text=BookInk,font=\figurefont\footnotesize] at (124,9)
      {$\lambda[E_g(\ell)-E_g(\ell-1)]$};
\end{tikzpicture}

 \caption{互斥类别条件下，最小费用流校准的一条可用路径。所有边容量为1，未标注费用的边费用为零；完整图包含全部合法槽位—物品边，以及每个类别的多个计数通道$\ell=1,\ldots,k$。物品入口到出口的共享容量限制重复，类别通道费用之和对应该类别最终计数的惩罚。此图依据正文目标绘制，省略未选路径。}
 \label{fig:rs-calibration-flow}
\end{figure*}

用六个候选、四个槽位核算。三个爱情片相关性为$(0.95,0.90,0.80)$，三个动作片为$(0.85,0.75,0.60)$；暂设各槽位价值相同，$p=(0.8,0.2)$、$\lambda=0.5$。爱情类前三条计数边费用约为$-0.1454,0.0279,0.0933$，动作类约为$0.0279,0.2012,0.2666$。负的第一条费用合法，它表示从零开始加入该类时的目标变化。

三种可行构成“爱情／动作”为$(3,1),(2,2),(1,3)$，各自最优相关性总和为3.50、3.45、3.15，KL散度约为0.0074、0.2231、0.7005，对应最大化形式的目标为1.7463、1.6134、1.2247。故选三个爱情片和最高分动作片。这个简单例子中，按同一增量目标逐项选择也可得到相同集合；最小费用流的区别在于它能够对整个可行图求最优，而非保证每个例子都胜过贪心。

当位置价值不同，局部选择就可能妨碍后续槽位。例如两件同类物品在两个剩余位置的价值分别为$(1,0.99)$与$(0.98,0)$；先把第一件放首位只得1，把第二件放首位、第一件放次位则得1.97，类别费用完全相同。整体流量求解能处理这种互斥分配。

一次流路由只能经过一个类别节点。多标签物品若要同时对所有标签计数，不能直接把其出口接到多个类别后仍宣称该图等价：流会选择其中一条，而不是复制全部贡献。应重新定义类别归属或验证新的建图。最优性针对具体图目标、固定候选和价值估计；位置边数随候选与列表长度的乘积增长，求解成本也要纳入服务预算。

% Outline: 5.2.2
\subsection{展示位置分配}
\label{sec:rs-position-allocation}

集合确定了对象身份，位置决定用户以什么次序、在什么区域遇到它们。若每件物品在每个位置的效用能够独立估计，令$v_{i,j}$为$i$在$j$的价值，沿用二值分配变量$M_{i,j}$，位置问题是
\begin{equation}
 \max_{\bm M}\sum_{i,j}M_{i,j}v_{i,j},
 \quad
 \sum_iM_{i,j}=1,\quad
 \sum_jM_{i,j}\leq1.
 \label{eq:rs-position-assignment}
\end{equation}
这是一个指派问题，可以用组合优化求解；资格禁止的对象—位置组合不应建立可用边。式中将收益相加，意味着其他对象的身份与位置不再改变$v_{i,j}$，这是重要假设。

例如，背包测评和水壶广告在“已查看”条件下的点击概率分别为0.30与0.10；教学假设首位、次位查看概率为0.8与0.4，并与对象独立。测评在前的预期总点击为$0.8\times0.30+0.4\times0.10=0.28$，广告在前则为0.20。这里乘法对应了查看与点击的条件概率链；如果所谓0.30本来已经包含位置查看，便不能再乘一次折扣。

若广告之间有最小间隔，或相邻对象会改变彼此的吸引力，式~\eqref{eq:rs-position-assignment}不足以表达目标。可以把邻接约束加入求解，也可以对完整排列估计效用；不能把不同页面条件下分别拟合的孤立收益直接相加。第~\ref{sec:rs-ads-allocation}~节的混排已决定插入位置，第~\ref{sec:rs-ad-load}~节若将位置并入行动，也已使用相应自由度；后续布局应尊重这些结果，或回到联合问题重新求解。

位置还包括二维布局和模块密度。同样两条广告，在四张首屏卡片中占一半，在十张卡片中只占五分之一；广告数量相同并不意味着可见负载相同。改变列数还会影响摘要可读性、触达顺序及对象大小。应记录真实区域、行列、滚动位置和可见范围，把页面执行结果与模型请求对应起来；位置纠偏见第~\ref{chap:rs-ranking}~章，跨请求的累计注意公平见第~\ref{chap:rs-constraints}~章。

% Outline: 5.2.3
\subsection{展示信息表达}
\label{sec:rs-presentation-expression}

用户看到商品标识并不足以判断是否合适，还需要知道价格、尺寸、防水等级或适用条件。展示表达决定用哪个标题、图片、摘要、商业标识和理由帮助用户判断。输入应包括可核验对象资料、当次需求与版面要求，输出则是一个有版本号的表达方案，便于把后续反馈关联到当时真实可见的信息。

同一款背包可以配两种教学摘要：“轻便防水，适合通勤”或“容量20升，售价399元；面料防泼水，拉链未作密封，不适合长时间淋雨”。如果资料支持后者，它帮助用户判断具体边界；前者把防泼水概括成防水，可能使读者形成过强预期。摘要更长也可能增加阅读成本，所以应同时检验事实完整性、理解和任务完成，而非只看点击率。图片和标题同样必须对应当前商品版本，商业标识也应在实际显示条件下检查。

语言模型可以生成表达候选，但生成依据与页面决策仍需分开。OneRec-Think在生成过程中使用的理由，不一定适合全部展示给用户；冗长推理会占空间，未经核验的属性也不能当成事实。广告创意及素材试验在第~\ref{chap:rs-advertising}~章展开。这里还须区分三个标准：描述是否真实，理由是否忠实于当前评分器，理由是否反映用户真实动机。三者可以分别成立或失败。

下面两种反事实方法针对第二个标准：固定模型，改变一部分输入，检查推荐决策是否改变。这里的“反事实”是模型输入的对照计算，不等于现实干预的因果识别，也不能代替第~\ref{sec:rs-uplift-ranking}~节的增量证据。即使一个理由能改变模型分数，仍需检查用户是否理解并认可这种表达。

% Outline: 5.2.3.1
\subsubsection{PRINCE：历史删除的反事实解释}
\label{sec:rs-prince}

某件背包有许多真实优点，但当前图推荐器可能只是因为用户以前购买过同类商品才把它排第一。\term{PRINCE}固定异质交互图与个性化PageRank（personalized PageRank, PPR）推荐器，寻找删除后使首选发生变化的尽量小的用户行为边集合\scite{rs-r112}PPR表示从用户节点反复随机行走、并以给定概率返回该用户后，各节点获得的稳态访问质量。

令$\bm W$为行随机转移矩阵，$\alpha\in(0,1)$为返回概率，则用户$u$的PPR行向量满足
\[
 \bm p_u=\alpha\bm e_u^\top+(1-\alpha)\bm p_u\bm W.
\]
在候选物品中选PPR最大的$a$。设$A_u$为可删除的用户出边，算法对替代候选$b$计算每条边$(u,n)$对$a$、$b$分差的贡献
\[
 d_n=W_{u,n}
 \left[\operatorname{PPR}_{\mathcal G\setminus A_u}(n,a)
      -\operatorname{PPR}_{\mathcal G\setminus A_u}(n,b)\right].
\]
这里的PPR在移除全部指定用户行为边的辅助图上预计算。论文的分解将删除子集后的分差符号化为剩余贡献之和；对固定$b$，按$d_n$从大到小移除，直到剩余和变为负，再在各替代对象之间取成功且边数较少的集合。

假设用户有两次购买边，通向$n_1,n_2$，权重各0.5。教学辅助图中，这两节点到$(a,b)$的PPR分别为$(0.30,0.10)$与$(0.05,0.10)$，于是贡献为0.10与$-0.025$，总和0.075支持$a$排在$b$前。删除第一条后只剩$-0.025$，比较方向反转；删除第二条则剩0.10，不能改变结果。输出可以表述为“在当前图模型中，不计入第一次购买后，$b$会超过$a$”，并附上被删除行为与重新核验的结果。

最小性依赖PPR分解、可删边范围及归一化条件；不应推广到任意神经评分器。若分差恰为零，还要使用实际并列规则核验，严格换序应要求替代分数更高。给定预计算PPR后，按替代物品逐一排序行为贡献的成本为$O(|C||A_u|\log|A_u|)$，另计PPR求解。删除操作改变的是固定模型输入，没有重新训练“真实偏好”；某条行为足以改变输出，也不说明它是用户现实购买的原因。

% Outline: 5.2.3.2
\subsubsection{CountER：属性变化的反事实解释}
\label{sec:rs-counter}

历史删除回答哪些行为支持选择，\term{CountER}考察当前对象的哪些属性支持它进入Top-$k$。先用交互监督训练一个读取用户和物品方面表示的评分器，解释时固定其参数，只优化物品属性的改变量\scite{rs-r113}设物品方面向量为$\bm z_i$，评分为$f_\theta(u,\bm z_i)$，当前第$k+1$项分数为$s_{\mathrm{cut}}$。原方法考虑降低属性表现，即$\bm\delta\leq\bm0$，希望用较少且较小的改动使目标掉出列表：
\begin{equation}
 \min_{\bm\delta\leq\bm0}
 \lVert\bm\delta\rVert_2+\gamma\lVert\bm\delta\rVert_0
 \quad\text{s.t.}\quad
 f_\theta(u,\bm z_i+\bm\delta)\leq s_{\mathrm{cut}}.
 \label{eq:rs-counter}
\end{equation}
$\lVert\bm\delta\rVert_0$计数被改动的属性。实际求解将它替为$\ell_1$正则，并将约束松弛为
\[
 \lVert\bm\delta\rVert_2+\gamma\lVert\bm\delta\rVert_1
 +\lambda\max\{0,\eta+f_\theta(u,\bm z_i+\bm\delta)-s_{\mathrm{cut}}\},
\]
其中$\eta>0$为换序裕量。优化后仍需重新计算排名，不能把松弛目标下降当成已经成功移出。

用餐厅“服务”和“价格实惠程度”两项作教学例，原值为$(4,3)$，评分器简化为$\sigma(0.6z_1+0.4z_2-2)$，原分数约0.832。若入选阈值为0.75，服务降低0.9得到$\sigma(1.06)\approx0.743$；实惠程度降低1.3得到$\sigma(1.08)\approx0.746$，两者都低于阈值。取原始计数代价$\gamma=0.1$，两个改动的复杂度分别为1.0与1.4，前者更简洁。这个比较只证明两种给定改动中前者代价小，没有证明找到了全部可行解中的最小改动。

生成解释时可说“若服务方面评分降低到3.1，当前模型就不会把它排入这份列表”，不能把修改后的3.1发布成餐厅事实。属性尺度、合法范围及现实可改变性也要检查。梯度求解次数决定解释成本，非凸模型和松弛约束可能导致失败或局部解；应报告成功率，而非只展示成功例子。模型忠实性还不能说明用户真实动机或现实中改善服务的因果收益。

% Outline: 5.2.4
\subsection{展示时间控制}
\label{sec:rs-presentation-time}

同一列表内的去重不能处理跨请求反复出现。时间控制需要定义何时允许曝光、两次曝光至少间隔多久，以及一个窗口内最多出现几次。统计对象可以是同一创意、同一活动、同一品牌或同一自然内容；窗口可以是固定日历日、滚动24小时或一次会话，口径改变就会改变计数。

例如，规则是“同一创意滚动24小时最多两次可见曝光，至少间隔30分钟”。教学请求在09:00产生第一次可见曝光，09:10被间隔规则排除，09:40可产生第二次；10:20被累计上限排除。被排除的请求不增加曝光计数。第二天09:01，第一条事件移出滚动窗口，若其他条件满足，可再次展示。应声明曝光是否要求实际可见，不能把服务器返回、卡片渲染和用户查看混成同一事件。

硬频控直接限制可行行动；软频次建模把重复状态纳入响应或价值预测，再由决策规则选择。前者能兑现明确上限，后者能表达用户、情境及对象之间的差异，但依赖足够的重复曝光样本。二者都需要可靠身份与窗口状态；跨设备关联缺失、计数延迟或并发请求会改变规则的实际执行。

当次重复价值还取决于需求是否变化。用户第一次只是浏览，第二次可能准备购买，也可能已经厌烦，因此不能先规定所有响应都随次数单调下降，再把它当成数据结论。下面的SFC学习这种条件关系，是否额外保留硬上限由展示规则决定。若要权衡现在展示对未来疲劳和可用机会的影响，还需第~\ref{chap:rs-long-term}~章的连续目标。

% Outline: 5.2.4.1
\Needspace{90mm}
\subsubsection{SFC：频次特征的点击预估}
\label{sec:rs-sfc}

\term{软频次控制}（soft frequency capping, SFC）在OFFSET点击模型中加入重复曝光特征，直接输出当前pCTR\scite{rs-m03}这里的名称不意味着模型本身选择曝光时机。设$f_{u,i,t}$是请求前、指定窗口及广告层级下的曝光次数，$B(f)$为分箱编号，$c(i)$指定共享权重组，则
\begin{equation}
 s'_{u,i,t}=s_{\mathrm{base}}(\bm x_t,i)+w_{c(i),B(f_{u,i,t})},
 \qquad
 \widehat p_{u,i,t}=\sigma(s'_{u,i,t}).
 \label{eq:rs-sfc}
\end{equation}
权重组可以全局共享，也可按活动或广告主区分。基础OFFSET利用用户、广告及情境特征的潜在因子估计响应；频次分箱权重与基础参数一起，以点击交叉熵和正则训练。

固定其他特征，取基础logit为$-2$，零次、1至2次、至少3次曝光的教学分箱权重分别为0、$-0.4$、$-0.9$。预测概率为0.1192、0.0832、0.0522。同一对象在不同重复状态下因此有不同价值。若第二个箱的一条成熟标签为未点击，忽略正则，交叉熵对该权重的梯度为$\widehat p-y=0.0832$，学习率0.1会把权重从$-0.4$更新至约$-0.4083$。

全局频次曲线有更多样本支持，却可能抹平活动差异；活动专属曲线更细，也更稀疏，需要正则、回退和足够覆盖。分箱不强制单调，允许某些重复状态具有较高响应。高频区域若被旧硬上限挡住，就缺少训练证据，不能仅靠模型外推认定解除上限有益。

输出pCTR可交给排序或广告价值预估。是否曝光、延后以及保留何种硬上限，应由相应决策规则确定。论文正式发表于CIKM 2019，2023是作者稿上传年份。点击率变化支持点击预测的评价，却不能单独量化主观厌烦、长期疲劳或留存。

% Outline: 5.2.5
\subsection{展示交互设计}
\label{sec:rs-presentation-interaction}

用户可以收藏、关闭、不感兴趣、筛选或修改要求，这些操作既改变当前界面，也为后续推荐提供证据。交互设计必须让入口文字与实际响应一致。“关闭这条广告”可以只移除当前卡片；“不再显示这个品牌”则要求后续候选按品牌范围处理。若用户关闭后立即看到同品牌另一条广告，是否违背预期取决于界面承诺，不能只靠内部把两条广告视为不同标识来解释。

应记录操作发生时的对象、表达版本、入口语义及可见位置。用户点击“不感兴趣”可能针对主题、商家、重复程度或当前时机；若需要区分，可提供成本适当的选项，但不能将未选择的原因补成标签。入口更容易发现会改变反馈发生率，不同界面下的关闭率必须连同机会和分母比较。

当系统还缺少需求信息时，询问本身也成为行动。第~\ref{chap:rs-understanding}~章说明了主动补充信息的价值，这里关心何时询问、何时推荐，以及拒绝后怎样继续。CRM学习询问与推荐时机，EAR在重复交互中利用拒绝更新，RecMind则用规划与工具访问组织执行。三者原任务和反馈来源不同，会话成功率、工具辅助预测与真实长期体验不能采用同一证据口径。

% Outline: 5.2.5.1
\subsubsection{CRM：询问与推荐的会话决策}
\label{sec:rs-crm}

把所有属性问完再推荐会增加操作成本，过早推荐又容易缺少关键条件。\term{CRM}用信念跟踪器、FM推荐器和策略网络共同决定下一步\scite{rs-r114}信念跟踪器为各属性输出可能取值的概率分布，保留“系统还有多不确定”；FM读取用户、物品及这些需求信息，产生候选评分；策略读取会话状态，在询问某个属性与触发推荐之间选择。

原文的信念跟踪器将对话文本的$n$元语法特征交给LSTM，再按属性分类，使用属性标注监督训练。FM按显式评分或相应交互目标训练，策略则通过会话奖励学习。设状态为$\bm s_t$，行动为$a_t$，从该轮起的折扣回报为$G_t$，策略梯度沿$G_t\nabla_\theta\log\pi_\theta(a_t\mid\bm s_t)$更新。成功、提问轮数和退出的奖励定义，决定模型如何权衡信息收益与交互成本。

“找酒吧”的教学会话开始时，城市信念为$(0.6,0.4)$。询问城市后，用户回答“北京”，对应信念更新为$(0.98,0.02)$，检索范围收缩，FM对两个合法候选给出4.2和3.6的评分。假设当前直接推荐的成功概率为0.60，询问后为0.80，成功奖励为1、每次提问成本为0.05，则一次询问的教学期望值为0.75，高于直接推荐的0.60；若再问一个属性只把成功率从0.80增至0.82，其值0.77便低于立即推荐的0.80。这是解释取舍的假设计算，真实概率与成本需证据支持。

输出是下一交互行动，触发推荐时才调用推荐器生成列表。原CRM围绕一次推荐触发的会话展开，多次拒绝后的持续推荐和模型更新由EAR扩展。属性空间不完整、语言理解错误或模拟用户过于配合，都可能改变策略效果。会话内的折扣回报也不自动涵盖数周后的留存。

% Outline: 5.2.5.2
\subsubsection{EAR：拒绝反馈下的交互更新}
\label{sec:rs-ear}

用户已经说明“喜欢摇滚”，仍可能拒绝当前推荐；这时重复同一列表无法利用新增信息。\term{估计—行动—反思}（estimation–action–reflection, EAR）让推荐器与会话策略交换更丰富的状态，并在拒绝之后更新推荐器\scite{rs-r115}设已确认属性集合为$P_u$，一种省去偏置和与候选无关项的FM评分写作
\[
 s(u,i,P_u)=\langle\bm h_u,\bm v_i\rangle
           +\sum_{p\in P_u}\langle\bm v_i,\bm e_p\rangle,
\]
其中$\bm e_p$为属性嵌入。离线训练既比较已交互与全局未交互物品，也比较与当前确认属性相容的候选内负例；另一个成对目标学习属性偏好，帮助策略决定询问什么。

策略状态连接四组信息：候选上的属性熵、模型预测的属性偏好、此前询问与推荐的结果，以及候选规模。属性熵高表示回答可能区分更多候选，偏好高表示用户可能认可该属性，两者不同。策略选择询问或推荐，再按成功、有效回答、退出及轮数等奖励训练。

若用户确认摇滚，候选从20个缩为8个；八个候选中四个现场、四个录音室版本，则“现场”这个二值属性的熵为$\log2$。即使模型预测用户偏好现场的概率不高，这个问题仍可能有较高区分能力。若策略推荐其中两个对象而被拒绝，EAR将这些真实展示对象作为比较负例，用历史已交互物品作为正例更新FM：
\[
 \mathcal L_{\mathrm{ref}}
 =-\sum_{i^+\in I_u^+,\,i^-\in I_{\mathrm{reject}}}
 \log\sigma\!\left[s(u,i^+,P_u)-s(u,i^-,P_u)\right]
 +\lambda\lVert\bm\theta\rVert_2^2.
\]
这里$I_u^+$是历史正例，$I_{\mathrm{reject}}$为本轮拒绝列表；在线更新不知道用户本轮真正想要的目标物品，不能把它偷偷当作训练正例。

教学分差若为$1.0-0.8=0.2$，数据损失对分差的导数为$-\sigma(-0.2)\approx-0.4502$，因此更新倾向于提高历史正例相对拒绝对象的分数。FM参数共享，其他候选也会受影响，更新后需重新评分并选择下一个问题或列表。拒绝可能只针对当次列表或版本，不等于永久不喜欢；过强更新会损害后续推荐。

论文主要在结构化属性反馈和用户模拟中检验这套交互循环，完整自然语言理解、生成错误与真实退出行为并未被同等覆盖。由模拟得到的更短轮数和较高成功率，应放在相应条件下解释；重复交互效果还需实际用户证据。

% Outline: 5.2.5.3
\subsubsection{RecMind：规划与工具调用的推荐执行}
\label{sec:rs-recmind}

语言模型的参数不包含最新目录和全部个人记录。\term{RecMind}在预训练语言模型外接入个人记忆、对象知识及SQL、检索、摘要等工具，让模型根据任务取得必要信息，再输出推荐或解释\scite{rs-r116}一次执行由规划、工具行动和工具观察组成；Self-Inspiring规划保留此前探索分支中的信息，使后续路径能够复用已经找到的证据，而非只保留最后选中的一条分支。

例如，用户要找500元以内的通勤背包。第一次只读查询返回历史中偏好的轻量和独立电脑隔层；第二次目录查询按价格、库存和隔层条件得到两件商品及其标识。某个探索分支查到其中一件仅防泼水，这条信息即使不属于最终主分支，也可以影响之后的摘要，避免输出“全天候防水”。最终对象、价格和属性应与工具结果逐一对应；如果工具没有返回防水信息，就不能靠语言模型补成肯定事实。

原框架用预训练LLM与提示组织推理，不要求为每个领域重新微调。它的五类离线任务包括评分预测、直接推荐、序列推荐、评论摘要和解释生成；这些任务考察如何使用工具完成既定输入，并不等于已经学习了CRM的主动询问策略，也不意味着可直接执行购买。

接入真实展示时，应明确工具能读哪些数据、输出能触发哪些界面动作、候选如何核验，以及失败时返回什么。规划步数、分支数量和工具延迟共同构成成本；工具结果过期或检索错误也会传递到最终回答。可读的文字不保证相关性、事实准确或评分器解释忠实，仍要分别检查这些标准。

% Outline: 5.3
\section{展示学习与评价}
\label{sec:rs-presentation-learning}

完整方案能够执行之后，还需要回答三个问题：已有反馈怎样更新决策依据，缺少的信息怎样取得，修改后的方案是否真的改善体验。学习、探索和评价可以共享日志，但职责不同。预测关闭概率不等于决定关闭代价，主动测试某种摘要也不等于已经证明它更好。

% Outline: 5.3.1
\subsection{展示反馈学习}
\label{sec:rs-display-feedback}

训练记录应还原实际执行的方案：展示对象及次序、位置和可见范围、表达版本、重复状态、交互入口，以及结果发生时间。仅保存排序器输出会丢掉执行过程中的过滤、补位和渲染失败，使标签被错误地配给未展示方案。跨请求训练还应只使用当时已可用的历史，防止把后续关闭或购买提前放入特征。

信号含义也随观察机会改变。假设同一广告在两个教学位置各被返回1000次，首屏有800次达到可见标准、产生16次关闭，下屏有200次可见、产生8次关闭。按返回次数计算，关闭率分别为1.6\%和0.8\%；按可见曝光计算，则为2\%和4\%。两者回答不同问题，且用户和可见条件仍可能不同，不能仅凭任一比率推断换位的因果效果。

若目标是“可见曝光后本次会话内是否关闭”，标签只有在会话结束并确认事件完整后才成熟；尚未结束的记录应等待或按明确的删失方法处理。未关闭并不表示满意，可能是入口难找；停留很长也可能是页面操作困难。点击、关闭、投诉、收藏和问卷应分别定义，再决定哪些预测进入内容、时间和交互策略。旧策略造成的选择偏差由第~\ref{chap:rs-ranking}~章和第~\ref{chap:rs-evaluation}~章承接。

\Needspace{85mm}
% Outline: 5.3.1.1
\subsubsection{Ad Close Mitigation：关闭反馈的分配修正}
\label{sec:rs-ad-close}

对于点击或收入相近的广告，关闭概率能够提供另一种区分依据。\term{Ad Close Mitigation}使用OFFSET特征因子模型预测广告关闭，并将预期关闭损失加入分配评分\scite{rs-m04}关闭事件作为正标签，成熟且未发生关闭的记录作为零标签，以用户、广告和情境特征训练二值概率$\widehat p_{\mathrm{close}}(\bm x,i)$；资格和曝光选择范围仍由原日志限定。

令$b_i$为每点击出价，$\widehat p_{\mathrm{click},i}$为点击概率，$c_{\mathrm{close}}$为一次关闭的指定代价，修正分数为
\begin{equation}
 s_i=b_i\widehat p_{\mathrm{click},i}
     -c_{\mathrm{close}}\widehat p_{\mathrm{close},i}.
 \label{eq:rs-ad-close}
\end{equation}
第一项是分配的出价加权响应，不能直接等同实际到账收入；第二项是按给定代价计算的预期损失。关闭概率由数据学习，代价系数表达平台取舍，允许的收入下降则是独立的评价约束。三者不能互相代替。

两个教学广告的第一项分别为0.10元和0.095元，关闭概率分别为0.03和0.005；取关闭代价1元，修正后分数为0.07元与0.09元，第二件因此超过第一件。若代价取零，结果回到原顺序。选择代价时应在独立验证和在线试验中检查关闭及真实收入，不能只用训练损失决定。

论文还相应调整含隐性成本的竞价与支付规则，分配改变后沿用旧收费公式未必一致，相关关系见第~\ref{chap:rs-advertising}~章。原文正式发表于WSDM 2020；机构网页的较早日期不代表会议年份。减少关闭只是改善一个可观察信号，不能据此断言所有打扰感或长期流失都已降低。

% Outline: 5.3.2
\subsection{探索曝光与信息获取}
\label{sec:rs-exploration-exposure}

没有展示过的新对象没有真实反馈，旧日志也可能只覆盖某些位置、摘要和频次。探索曝光主动选择一部分未知方案，付出当前机会成本来取得信息。探索单位应与问题一致：测试新对象就保持表达规则稳定，测试摘要就固定对象与位置；若同时更改多个变量，则应评价整个联合方案，或设计能够分辨交互的试验。

随机试验为比较当前方案提供识别依据，探索策略则可能为了改进未来决策主动偏向信息不足的行动。它们可以结合，但不能混称。候选资格、用户控制和测试预算先限定允许行动，再由策略选择；记录所执行的行动、选择机制及必要概率，后续才有可用日志。

预测不确定性、行动概率与策略价值估计的置信区间也各不相同。上置信界（upper confidence bound, UCB）把预测均值与不确定性奖励相加，用于在利用和探索之间选择；若相同状态下始终选最大值，它就是确定性行动，不会自动为其他行动提供随机试验支撑。下面分别考虑只观察点击前缀的列表接口，以及单个神经上下文行动接口。

% Outline: 5.3.2.1
\subsubsection{CascadeLinUCB与CascadeLinTS：部分点击反馈的列表探索}
\label{sec:rs-cascade-linear}

一组推荐中，用户可能从上到下查看，遇到第一个有吸引力的对象就点击并停止。\term{级联模型}（cascade model）明确采用这一检查规则；首点击之前的未点击对象得到失败标签，首点击得到成功标签，后面的对象没有被检查。若全列表无点击，则在该模型下视为全部检查且均无吸引。\term{CascadeLinUCB}与\term{CascadeLinTS}进一步用对象特征共享吸引概率的统计证据\scite{rs-r117}

设对象$i$特征为$\bm x_i\in\mathbb R^d$，真实吸引概率为$w_i=\langle\bm x_i,\bm\theta^\star\rangle\in[0,1]$，并假设各对象吸引事件独立。列表至少得到一次点击的概率是
\begin{equation}
 U(L)=1-\prod_{i\in L}(1-w_i).
 \label{eq:rs-cascade-utility}
\end{equation}
这个目标只取决于所选集合，不取决于顺序；但顺序影响哪些对象得到观察，因而影响学习过程。已知$w_i$时选择最大的$k$个即可，未知时则需探索。

为简化记号，令噪声缩放参数为1。对所有已经实际检查的对象，累计
\[
 \bm A=\bm I+\sum_{\mathrm{observed}}\bm x_i\bm x_i^\top,
 \qquad
 \bm b=\sum_{\mathrm{observed}}y_i\bm x_i,
 \qquad
 \widehat{\bm\theta}=\bm A^{-1}\bm b.
\]
CascadeLinUCB按$\min\{1,\bm x_i^\top\widehat{\bm\theta}
+c\sqrt{\bm x_i^\top\bm A^{-1}\bm x_i}\}$选择前$k$项，其中$c$按置信要求及模型条件确定。CascadeLinTS则从$\mathcal N(\widehat{\bm\theta},\bm A^{-1})$采样一组参数，再按$\bm x_i^\top\widetilde{\bm\theta}$选择；这个高斯分布是算法的参数采样近似，不是伯努利似然下自动精确的后验。

三个教学候选$a,b,c$的均值估计依次为0.45、0.40、0.30，对应UCB奖励为0.02、0.02、0.20。相加后$c$与$a$最高，故两项列表可取$(c,a)$，让不确定的新对象$c$获得机会。若首位$c$被点击，只更新$(c,1)$，不能给$a$记零；若第二位$a$被点击，更新$(c,0),(a,1)$；若无点击，才更新两个零。特征更新还会影响与它们相似的其他候选，这使学习可以在对象间泛化。

图~\ref{fig:rs-cascade-feedback}把这三种可见前缀画出，关键是区分“未点击”与“未观察”。若实际用户会多次点击、跳着查看或被相邻对象改变偏好，原反馈假设就不成立。线性泛化与置信界也有条件：原文对相应假设下的UCB给出遗憾界，对TS主要提供经验检验，不能把同一理论保证同时赋予两者。它们优化给定吸引模型下的累积回报，并未建模长期兴趣被曝光改变。

\begin{figure}[htbp]
 \centering
 % 固定列表 (c,a)，三种结果只更新被级联检查覆盖的前缀。
\begin{tikzpicture}[
 x=1mm,y=1mm,font=\figurefont\small,
 checked/.style={draw=BookBlue,fill=BookBlue!6,line width=0.85pt,
                 minimum width=21mm,minimum height=9mm},
 unseen/.style={draw=BookMuted,dashed,line width=0.85pt,
                minimum width=21mm,minimum height=9mm,text=BookMuted}
]
 \path[use as bounding box] (0,0) rectangle (112,47);
 \node[font=\figurefont\footnotesize,text=BookMuted] at (40,44) {位置1};
 \node[font=\figurefont\footnotesize,text=BookMuted] at (68,44) {位置2};
 \node[font=\figurefont\footnotesize,text=BookMuted,anchor=west]
      at (84,44) {用于更新};
 \node[anchor=west] at (0,34) {首位点击};
 \node[checked] (r11) at (40,34) {$c:1$};
 \node[unseen] (r12) at (68,34) {$a:\ ?$};
 \node[anchor=west] at (84,34) {$c$};
 \node[anchor=west] at (0,20) {次位点击};
 \node[checked] (r21) at (40,20) {$c:0$};
 \node[checked] (r22) at (68,20) {$a:1$};
 \node[anchor=west] at (84,20) {$c,a$};
 \node[anchor=west] at (0,6) {没有点击};
 \node[checked] (r31) at (40,6) {$c:0$};
 \node[checked] (r32) at (68,6) {$a:0$};
 \node[anchor=west] at (84,6) {$c,a$};
 \draw[->,BookMuted,line width=0.85pt] (r21.east) -- (r22.west);
 \draw[->,BookMuted,line width=0.85pt] (r31.east) -- (r32.west);
\end{tikzpicture}

 \caption{两对象级联列表的可观察前缀。实线框为已检查对象，虚线框为未观察对象；只有已检查项进入统计更新。无点击时两项均记零依赖完整级联检查假设，不适用于用户可能中途退出而未检查剩余项的情形。}
 \label{fig:rs-cascade-feedback}
\end{figure}

% Outline: 5.3.2.2
\subsubsection{NeuralUCB：神经响应模型的置信探索}
\label{sec:rs-neuralucb}

固定线性特征可能不足以描述用户和对象的响应关系。\term{NeuralUCB}用神经网络$f_{\bm\theta}(\bm x_{t,a})$估计行动$a$的奖励均值，再利用网络参数梯度衡量该方向积累了多少证据\scite{rs-r118}设网络宽度为$m$，归一化梯度为$\bm g_{t,a}=\nabla_{\bm\theta}f_{\bm\theta}(\bm x_{t,a})/\sqrt m$，累计矩阵初值为$\bm Z^{(0)}=\lambda\bm I$，选择分数为
\begin{equation}
 B_{t,a}=f_{\bm\theta^{(t-1)}}(\bm x_{t,a})
          +\beta_t\sqrt{\bm g_{t,a}^\top
                      (\bm Z^{(t-1)})^{-1}\bm g_{t,a}},
 \quad a_t=\argmax_a B_{t,a}.
 \label{eq:rs-neuralucb}
\end{equation}
所有候选的均值和梯度在本轮反馈前计算。原算法的置信系数依赖噪声、网络、训练精度和置信要求，$\beta_t$不是任意网络都适用的固定常数。

得到所选行动的奖励$r_t$后，先用该轮选择时的梯度更新$\bm Z^{(t)}=\bm Z^{(t-1)}+\bm g_{t,a_t}\bm g_{t,a_t}^\top$，再以已选择样本训练网络。原文训练目标包含围绕初始参数$\bm\theta^{(0)}$的正则：
\begin{equation}
 \mathcal L_t(\bm\theta)=
 \frac12\sum_{s=1}^{t}
       [f_{\bm\theta}(\bm x_{s,a_s})-r_s]^2
 +\frac{m\lambda}{2}
       \lVert\bm\theta-\bm\theta^{(0)}\rVert_2^2.
 \label{eq:rs-neuralucb-train}
\end{equation}
梯度统计与网络训练是相互配合的两部分；只给普通CTR网络加一个未经定义的方差项，不等于实现了NeuralUCB。

在一个推荐测试位上，教学候选$a,b$的预测均值为0.40、0.35。假设归一化梯度为$(1,0)$与$(0,1)$，$\bm Z=\operatorname{diag}(100,4)$，取$\beta_t=0.2$，置信奖励为0.02与0.10，分数变为0.42与0.45，因而选择$b$。完成反馈后，矩阵第二个对角元由4增至5；暂固定梯度，下一次这一方向的奖励降为$0.2/\sqrt5\approx0.0894$。这个计算说明了信息积累怎样减少探索奖励，实际网络更新还会改变均值和梯度。

NeuralUCB输出本轮行动，不能直接被解释为完整列表策略；多位置还需定义组合行动和可观察反馈。全参数梯度矩阵的大小随参数数平方增长，实际近似也应说明与原算法的差别。理论结论依赖足够网络宽度、奖励噪声、训练及上下文条件，不为任意漂移环境提供无条件保证。确定性最大化也不产生对所有行动的正概率覆盖。

% Outline: 5.3.3
\subsection{体验与方案评价}
\label{sec:rs-display-experience}

评价应以用户实际收到的方案为单位，而不是只评价一个中间分数。内容、位置、表达、时间和交互共同决定体验；“营销感”也由这些变量共同形成。广告是否符合当前需求、商业标识是否明确、首屏是否拥挤、同一品牌是否反复出现、关闭后是否得到可理解的响应，都可能参与其中，不能把广告数量当作全部解释。

单变量比较能回答局部问题。例如保持对象、位置、摘要与数量相同，只改变同一广告的最小重复间隔；或固定广告数量，比较不同位置与摘要的联合方案。后者若同时改善，不应把全部收益归给某一个变量。需要区分交互时，可在可执行范围内交叉安排位置与表达，再按第~\ref{chap:rs-evaluation}~章的在线协议比较。

结果至少覆盖当前任务与相关体验证据。任务完成、有效观看或购买是行为结果；关闭、投诉和拒绝是部分负反馈；问卷能够直接询问理解、控制感与打扰感，却受邀请和回答选择影响。点击升高可能来自信息更清楚，也可能来自夸张标题，应结合后续结果和理解检查判读。反事实理由则还应单独报告改变模型决策的成功率、改动规模及解释成本。

汇总平均值后，应检查不同用户、会话阶段和重复状态，并保持标签成熟及观察窗口一致。模型只在旧展示条件下预测准确，不代表改变页面仍会准确；探索收集到少量信息，也不等于效果估计已有足够精度。若目标转为长期疲劳、留存或内容生态，需要采用下一阶段连续决策与较长观察窗口，不能把本次页面的代理收益直接外推。

% Outline: 5.4
\section{本章小结}
\label{sec:rs-presentation-summary}

三个高分对象是否值得一起出现，取决于它们是否重复、互补，放在什么位置，用什么信息呈现，以及此前已经出现多少次。完整展示由内容、位置、表达、时间和交互共同定义。候选重排、目录列表生成、来源混排及会话策略各自控制其中一部分，只有输入范围、固定条件和输出接口都明确，方法才能组合和比较。

整组差异与兴趣构成匹配具有不同目标，真实描述、模型忠实解释与用户动机也有不同证据要求。关闭概率、频次响应、奖励模型和探索置信项能够支持选择，却不能替代实际方案的体验评价。日志要记录执行而非仅记录计划，未观察反馈也不能填成失败。

展示方案通常同时带来收益、体验变化和资源消耗。即使所有响应都已准确估计，仍需决定哪些量可以折算、哪些要求必须满足、哪些参与方的机会需要保护。下一章在本章的可执行方案上建立这些目标与约束，使“能展示什么”进一步变成“应按什么条件作出选择”。
